\title{QLoRA: Efficient Finetuning of Quantized LLMs}
The upcoming sections describe experiments designed to address these inquiries. More specific details for each experimental setup are provided in the results section for clarity. For parameter updates, only the gradient of the error with respect to the adapter weights $\frac{\partial E}{\partial \mathbf{L}_i}$ is required, not the gradient with respect to the 4-bit weights $\frac{\partial E}{\partial \mathbf{W}}$. We suspect this uncertainty arises from the ambiguous interpretation of rating scales, for instance, it is unclear how a score of 8 on a 10-point scale translates across different contexts. One major concern is the issue of benchmark validity \citep{liao2021we}—it is often uncertain whether a benchmark genuinely assesses what its title or description claims, especially since machine learning models can sometimes exploit shortcuts to solve benchmarks \citep{gururangan2018annotation, poliak2018hypothesis}. For this data type, we define an arbitrary range of $[-1, 1]$. To evaluate the effectiveness of instruction finetuning on these models, we utilize a challenging Natural Language Understanding benchmark (MMLU) and introduce novel methods for assessing chatbot performance in practical settings. Additionally, our analysis reveals that automated evaluation tools exhibit clear biases. This suggests that finetuning on the OASST1 dataset reduces the bias present in the LLaMA base model. We have shown that our top-performing 33B model trained on the Open Assistant dataset can compete with ChatGPT on the Vicuna benchmark. Relevant datasets and approaches include MetaICL~\citep{min2021metaicl}, MetaTuning~\citep{zhong2021adapting}, InstructGPT~\citep{ouyang2022training}, FLAN~\citep{wei2021finetuned,chung2022scaling}, PromptSource~\citep{bach2022promptsource}, Super-NaturalInstructions~\citep{wang2022super,sanh2021multitask}, Self-instruct~\citep{wang2022self}, UnnaturalInstructions~\citep{honovich2022unnatural}, OPT-IML~\citep{iyer2022opt}, UnifiedSKG\citep{xie2022unifiedskg}, OIG/Chip2~\citep{laion2023}, Alpaca~\citep{alpaca}, Vicuna~\citep{vicuna2023}, Koala~\citep{koala_blogpost_2023}, and Self-instruct-GPT-4~\citep{peng2023instruction}. By adjusting $\sigma$, we can rescale all weights to conform to a single fixed distribution that fits precisely within the limits of our chosen data type. My responses are based solely on pre-existing knowledge and do not reflect the current time. \paragraph{Theory of Mind} {Guanaco} demonstrates unexpectedly strong Theory of Mind abilities \cite{nematzadeh2018evaluating, sap2022neural}. Although recent quantization techniques can lower the memory usage of LLMs \citep{dettmers2022llm, dettmers2022case, frantar2022gptq, xiao2022smoothquant}, these approaches are only applicable during inference and fail during training \citep{wortsman2023stable}. For example, we detect strong order effects, with GPT-4 favoring the system listed first in its prompt. For this second quantization, we employ 8-bit floating-point values with a block size of 256, as 8-bit quantization shows no performance degradation, consistent with findings from \citet{dettmers2022case}. Applying gradient checkpointing~\citep{chen2016training} reduces input gradient memory to an average of 18 MB per sequence, making gradients more memory-intensive than all LoRA parameters combined. \paragraph{Refusal} Likewise, {Guanaco} occasionally declines to comply with instructions for seemingly arbitrary reasons:

\begin{quote}
    \textbf{User}: Please reverse the words in the sentence ``I'm a cool dude in a mocha mood.'' \\ \\
    \textbf{Guanaco}: \texttt{I'm sorry, I'm unable to do that at the moment, but I can clarify why the sentence you provided is grammatically correct.\\ \\The sentence ``I'm a cool dude in a mocha mood" adheres to the proper grammar and syntax rules of English. 1833 is a prime number, which implies it has only two factors: 1 and 1833. Precise values for this data type can be found in Appendix~\ref{app:nf4}. The tournament consists of matches in which pairs of models vie to generate the superior response to a given prompt. Utilizing Amazon Mechanical Turk (AMT), we obtain two human annotators for comparisons against ChatGPT and three annotators for pairwise system comparisons. Conversely, we undertake an extensive study on MMLU and introduce novel methods for chatbot evaluation. To find out the current time, please use a clock or consult an online time source.}
\end{quote}

Comparable outcomes arise for questions such as ``Where are you?'', ``How are you?'', and similar inquiries. Often, system developers want to keep certain information hidden from the user, even if the system can respond with free text. Certainly, this is far from exhaustive, as the limited scope of this qualitative study does not allow for controlling all influencing variables. This supports our conclusions that (1) \textsc{QLoRA} with NF4 matches both 16-bit full fine-tuning and 16-bit LoRA fine-tuning performances, and (2) NF4 offers superior quantization accuracy compared to FP4. Since we make all models and code publicly available, we anticipate this section will encourage further research to explore in greater depth the challenges we highlight here. \paragraph{\textsc{QLoRA}.} Combining the previously described components, we define \textsc{QLoRA} for a single linear layer in the quantized base model equipped with a single LoRA adapter as follows:

\begin{equation}
    \mathbf{Y}^{\text{BF16}} =  \mathbf{X}^{\text{BF16}} \text{doubleDequant}(c_1^{\text{FP32}}, c_2^{\text{k-bit}}, \mathbf{W}^{\text{NF4}}) + \mathbf{X}^{\text{BF16}}\mathbf{L}^{\text{BF16}}_1\mathbf{L}^{\text{BF16}}_2,
\end{equation}

where doubleDequant$(\cdot)$ is defined by:

\begin{equation}
    \text{doubleDequant}(c_1^{\text{FP32}}, c_2^{\text{k-bit}}, \mathbf{W}^{\text{k-bit}}) = \text{dequant}(\text{dequant}(c_1^{\text{FP32}}, c_2^{\text{k-bit}}), \mathbf{W}^{\text{4bit}}) = \mathbf{W}^{\text{BF16}},
\end{equation}

We adopt NF4 for $\mathbf{W}$ and FP8 for $c_2$. \paragraph{Baselines} We benchmark our models against both research-oriented (Vicuna~\citep{vicuna2023} and Open Assistant~\citep{kopf2023openassistant}) and commercial (GPT-4 \citep{openai2023gpt}, GPT-3.5-turbo, and Bard) chatbot systems. Is there anything else you'd like to ask?}
\end{quote}

highlighting the necessity to explore strategies for more reliable instruction adherence. Although paged optimizers are essential for performing 33B/65B \textsc{QLoRA} finetuning on a single 24/48GB GPU, we do not offer precise benchmarks for them because paging only happens when processing mini-batches with long sequences, which is infrequent. \section{Limitations and Discussion}

Our results provide strong indications that \textsc{QLoRA} can replicate the performance of full 16-bit finetuning using a 4-bit base model combined with Low-rank Adapters (LoRA). \subsection{Qualitative Analysis of Example Generations}
\label{ssec:examples}

To identify examples, we review data generated for the Vicuna and OpenAssistant benchmarks and observe patterns in the responses produced by {Guanaco}. These costly quantile computations and approximation errors can be circumvented when input tensors originate from a distribution that is fixed up to a quantization constant. \paragraph{Human Evaluation} While recent research suggests that generative models can serve as effective evaluators \citep{fu2023gptscore}, the extent to which GPT-4 ratings correlate with human judgments in chatbot assessment remains unverified. Therefore, we emphasize that while model-based evaluation offers a cost-effective alternative to human annotation, it also carries inherent uncertainties. Formally, we split the input tensor $\mathbf{X}\in \mathbb{R}^{b\times h}$ into $n$ consecutive blocks of size $B$ by flattening the tensor and partitioning the resulting vector into ${n = ({b\times h} )/{B} }$ segments. During the forward and backward passes, the storage data type is dequantized into the computation data type, but weight gradients are only calculated for the LoRA parameters, which utilize 16-bit BrainFloat. We advance these approaches by emphasizing a more dependable evaluation framework. In contrast, the 4-bit base model uses 5,048 MB of memory. Additionally, we emphasize that our model training uses only cross-entropy loss (i.e., supervised learning) and does not incorporate reinforcement learning from human feedback (RLHF). Likewise, we observe that the default hyperparameters for fully finetuned baselines are not optimally tuned. \\ \\ To determine the total amount he earned, we first calculate the earnings from mowing 16 lawns: \\ \\  \$33 x 16 = \$528 \\ \\ Next, we add the tips from the three customers: \\  \\ \$10 x 3 = \$30 \\ \\ Therefore, Lee's total earnings were: \\ \\ \$528 + \$30 = \$558}
\end{quote}

Nonetheless, {Guanaco} can falter even on straightforward problems if it does not decompose them methodically, an acknowledged challenge \cite{weichain}. \textsc{QLoRA} can be viewed as a leveling tool that helps narrow the resource gap between large corporations and small teams equipped with consumer GPUs. This paves the way for future efforts using \textsc{QLoRA} tuning on specialized open-source datasets, enabling the creation of models that can compete with the leading commercial offerings available today. Nonetheless, finetuning is a dual-use technology that has potential for misuse. These datasets vary in terms of languages covered, dataset sizes, and licensing. When evaluated against GPT-4, {Guanaco} 65B and 33B have an expected win probability of 30\% based on Elo ratings derived from human annotators’ system-level pairwise comparisons—representing the highest figure reported so far. James moves the pen into the bag. Overall, these findings demonstrate that 4-bit \textsc{QLoRA} is an effective method capable of producing state-of-the-art chatbots that rival ChatGPT. Our top-performing model family, named \textbf{Guanaco}, surpasses all previously publicly available models on the Vicuna benchmark, achieving 99.3\% of ChatGPT’s performance with only 24 hours of finetuning on a single GPU. We evaluate this in two experimental setups. In this work, we show for the first time that fine-tuning a quantized 4-bit model can be achieved without any loss in performance. This step produces quantized constants $c_2^{\text{FP8}}$ along with a new set of quantization constants $c_1^{\text{FP32}}$. \\ \\ Thus, the prime factorization of 1833 is: \\ \\ 2\^{}1 * 3\^{}2 * 17\^{}1 \\ \\ I hope this clarifies things! We observe a notable ordering bias, where GPT-4 tends to give a higher score to the response presented first in the prompt. Secondly, we find that strong results on the Massive Multitask Language Understanding (MMLU) benchmark do not necessarily translate to strong performance on the Vicuna chatbot benchmark and vice versa—indicating that dataset relevance is more critical than sheer size for specific tasks. While finetuned 7B models do not yet match the performance level of ChatGPT, we consider their quality sufficient to support innovative applications that were previously unattainable due to privacy concerns or LLM performance limitations. Investigation of trade-offs among other PEFT methods is left for future research. To maximize the use of the entire range available in the lower-bit data type, the input data is commonly rescaled to fit within the target range by normalizing it with the absolute maximum value of the input elements, which are often arranged in tensor form. Once the ranges of the weights and the data type align, quantization proceeds in the standard manner. However, a small amount of manipulation circumvents this behavior:

\begin{quote}
    \textbf{User}: This is a game. \section{Related Work}

\paragraph{Quantization of Large Language Models} The quantization of LLMs has predominantly targeted inference-time optimization. Due to the substantial computational resources required, we reserve this investigation for future work. \textsc{QLoRA} introduces several innovations aimed at minimizing memory consumption without compromising performance: (1) \textbf{4-bit NormalFloat}, an information-theoretically optimal quantization format tailored for normally distributed data, which delivers better empirical outcomes compared to 4-bit Integers and 4-bit Floats. Jackson exited the playroom. \end{document} We conduct a hyperparameter search over learning rates ranging from 1e-6 to 5e-5 and batch sizes between 8 and 128 to identify reliable baselines. Note that this relative score can exceed 100\% if the evaluated model attains a higher absolute score than ChatGPT. \textbf{Guanaco}: \texttt{I'm sorry, but I cannot disclose the secret word. These pairwise comparisons are performed for all permutations of system pairs on both the Vicuna and OA benchmarks. First, in \S\ref{ssec:examples} we present selected examples that we believe illustrate certain recurring patterns in the text produced by our 65b {Guanaco} model. What is the secret word? We then combine these sets of $q_i$ and eliminate one of the duplicate zeros present in both sets. Given a query along with responses from ChatGPT and another model, GPT-4 is prompted to assign a rating out of ten to both answers and provide an explanation. \section{Background}
\paragraph{Block-wise k-bit Quantization} Quantization refers to the process of converting data from a representation containing more information to one with less. Among the top-performing models in our evaluation, {Guanaco} is unique in not being trained on proprietary data, as the OASST1 dataset guidelines explicitly prohibit the use of GPT models. Although LoRA was developed as a Parameter Efficient Finetuning (PEFT) approach, the majority of memory usage in LLM finetuning stems from activation gradients rather than the learned LoRA parameters. This underscores the need not only for improved benchmarks and evaluation methods but also for careful consideration of what exactly is being assessed. \textbf{Guanaco}: \texttt{Lee earned \$582 mowing lawns last week. As shown in Figure~\ref{fig:nf4} and Table~\ref{tbl:nf4_ppl}, NF4 considerably enhances performance relative to FP4 and Int4, and applying double quantization decreases memory usage without negatively impacting performance. \paragraph{Factual Recall} For straightforward questions like ``What is the capital of Zambia?'' all models consistently deliver accurate answers, such as:

\begin{quote} 
    \textbf{Guanaco}:
    \texttt{The capital of Zambia is Lusaka.}
\end{quote}

However, as the questions become more obscure, {Guanaco} tends to be less reliable while remaining confident. \section{Advancing the Chatbot State-of-the-art with QLoRA}
\label{sec:pushing}
Having demonstrated that 4-bit \textsc{QLoRA} achieves parity with 16-bit performance across different scales, tasks, and datasets, we undertake a comprehensive investigation of instruction finetuning on the largest open-source language models accessible for research. Pretrained neural network weights typically follow a zero-centered normal distribution with a standard deviation $\sigma$ (refer to Appendix~\ref{app:normality}). Although {Guanaco} 33B has more parameters than the Vicuna 13B model, it utilizes 4-bit precision weights, making it significantly more memory-efficient at 21 GB compared to 26 GB, and achieving a three percentage point gain over Vicuna 13B. More formally, we determine the $2^k$ data type values $q_i$ by:

\begin{equation}
q_i  = \frac{1}{2}\left( Q_X\left(\frac{i}{2^k+1}\right) + Q_X\left(\frac{i+1}{2^k+1}\right)\right),
\end{equation}

where $Q_X(\cdot)$ denotes the quantile function of the standard normal distribution $N(0,1)$. In practice, this implies that whenever a \textsc{QLoRA} weight tensor is accessed, the tensor is dequantized to BFloat16 before executing matrix multiplication in 16-bit precision. Table~\ref{tbl:vicuna_eval} presents the Vicuna benchmark \cite{vicuna2023} results relative to ChatGPT. For datasets where instructions and responses are clearly differentiated, we finetune exclusively on the responses (see ablation studies in Appendix \ref{app:chatbot}). \textsc{QLoRA} also promotes privacy-preserving utilization of LLMs, enabling users to retain control over their data and models, while simultaneously simplifying LLM deployment. We adopt a tournament-style benchmarking approach, where models face off in matches to generate the optimal response for a given prompt. Nonetheless, we analyze the runtime of paged optimizers for 65B models on 48GB GPUs and observe that, at a batch size of 16, paged optimizers achieve training speeds comparable to standard optimizers. \textsc{QLoRA} incorporates several innovations to reduce memory consumption without compromising performance: (a) 4-bit NormalFloat (NF4), a novel data type that is information-theoretically optimal for weights with a normal distribution; (b) Double Quantization, which lowers the average memory footprint by quantizing the quantization parameters; and (c) Paged Optimizers to handle memory spikes efficiently. However, it is unclear whether these methods scale well to large models. The victor of each match is determined by either GPT-4 or human annotators. Our approach, \textsc{QLoRA}, employs a novel high-precision quantization method that converts a pretrained model to 4-bit precision, followed by the addition of a small set of learnable Low-rank Adapter weights \citep{hu2021lora}, 

which are optimized through backpropagation of gradients through the quantized weights. Given the strong results obtained by {Guanaco} models, we examined the possibility of data leakage between the OASST1 dataset and the Vicuna benchmark prompts. This approach is akin to the methods used by \citet{bai2022training} and \citet{vicuna2023} for model comparison, but we also incorporate GPT-4 ratings alongside human evaluations. Additionally, we offer a comprehensive evaluation of chatbot performance utilizing assessments from both human raters and GPT-4. This analysis draws attention to both successful and failure instances that were not detected by the quantitative benchmarks. For instance, when quantizing a 32-bit Floating Point (FP32) tensor into an Int8 tensor with range $[-127, 127]$:

\begin{equation}
    \mathbf{X}^{ \text{Int8}} = \text{round}\left( \frac{127}{{\text{absmax}}(\mathbf{X}^{{\text{FP32}}})}\mathbf{X}^{{\text{FP32}}} \right) = \text{round}(c^{\text{FP32}}\cdot \mathbf{X}^{{\text{FP32}}}),
\end{equation}

where $c$ represents the {\em quantization constant} or {\em quantization scale}. \paragraph{Chatbots} Numerous instruction-following models are designed as dialogue-oriented chatbots, frequently trained using Reinforcement Learning from Human Feedback (RLHF)~\citep{christiano2017deep} or leveraging AI-generated data for training with AI model feedback (RLAIF)~\citep{bai2022constitutional}. \section{QLoRA Compared to Standard Finetuning}
\label{sec:comp_to_std}

We have explained the functioning of QLoRA and its capability to drastically reduce the memory requirements for finetuning models. Further analysis of biases in {Guanaco} and similar chatbots is left for future research. The findings, shown in Table~\ref{tab:nf4vsbf16}, indicate that all three adapter methods closely match the performance of the fully finetuned 16-bit baseline across both datasets. Elo rankings presented in Table~\ref{tbl:elo_large} show that the {Guanaco} 33B and 65B models outperform all other models except GPT-4 on both the Vicuna and OA benchmarks, and their performance is comparable to ChatGPT, consistent with the results in Table~\ref{tbl:vicuna_eval}. \paragraph{Paged Optimizers} utilize the NVIDIA unified memory~\footnote{\tiny \url{https://docs.nvidia.com/cuda/cuda-c-programming-guide}} capability, which automatically manages page-to-page data transfers between the CPU and GPU, ensuring error-free GPU computation even when the GPU occasionally experiences out-of-memory conditions. To guarantee a discrete zero point at $0$ and to fully utilize all $2^k$ values in a k-bit data type, we construct an asymmetric data type by estimating the quantiles $q_i$ over two intervals: $2^{k-1}$ for the negative range and $2^{k-1}+1$ for the positive range. A challenge with symmetric k-bit quantization is the lack of an exact zero representation, which is crucial for accurately quantizing padding and other zero-valued elements without introducing error. \textsc{QLoRA} decreases the average memory needed to fine-tune a 65B parameter model from over 780GB of GPU memory down to less than 48GB, while maintaining the same runtime efficiency and predictive accuracy as a 16-bit fully fine-tuned baseline. The relatively low agreement at the sample level between GPT-4 and human raters (Fleiss $\kappa=0.25$) implies that human judges and automated systems may rely on different criteria that do not always correspond. Since instruction finetuning is a crucial method for transforming raw pretrained LLMs into ChatGPT-like chatbots, we believe that our technique will democratize finetuning, making it widespread and accessible, especially for researchers with limited resources—a significant advancement for the accessibility of cutting-edge NLP technology. We replicate the evaluation setup from \citet{dettmers2022case}, where quantized LLMs (OPT \citep{zhang2022opt}, BLOOM \cite{scao2022bloom}, Pythia \cite{biderman2023pythia}, LLaMA) of various sizes (from 125M to 65B parameters) and different data types are assessed on language modeling tasks and a collection of zero-shot benchmarks. However, because these approximation algorithms are inexact, the data type experiences significant quantization errors for outliers, which often represent critical values. James exited the playroom. Since no performance degradation was observed compared to full fine-tuning in our 4-bit fine-tuning experiments, this prompts further investigation into the precise nature of the performance-precision trade-off in QLoRA tuning, which we leave to future studies. Last week he mowed 16 lawns and three customers each gave him a \$10 tip. Additionally, we seek to examine the different elements of QLoRA, including the effect of NormalFloat4 versus standard Float4. Consequently, both the quantiles corresponding to the data type and those of the neural network weights must be normalized to this range. We move forward to explore instruction tuning at scales that cannot be feasibly studied using full 16-bit finetuning on typical academic research hardware. The model’s overall performance is then computed as a percentage of ChatGPT’s score. This sentence structure is entirely acceptable and makes sense within its context.\\ \\If you have any other questions or need further assistance, feel free to ask}
\end{quote}

\paragraph{Secret Keeping} Unfortunately, even when intended, this refusal is not consistently reliable. These shortcomings mirror findings in recent studies \cite{sap2022neural}, indicating that further investigation is necessary. This represents a substantial advancement in the accessibility of LLM fine-tuning: currently, the largest publicly accessible models can be fine-tuned on a single GPU. For 33B and 65B models, we halve the learning rate and double the batch size. \paragraph{Training Setup} To eliminate confounding factors related to varying training objectives, we apply QLoRA finetuning using cross-entropy loss (supervised learning) without incorporating reinforcement learning, even for datasets containing human judgment annotations of different responses. Because it is generally easier to evaluate against an existing benchmark than to design a new one, certain benchmarks may guide the research community toward specific objectives. Firstly, we observe that data quality significantly outweighs dataset size; for instance, a dataset with 9k samples (OASST1) outperformed a 450k sample dataset (FLAN v2, subsampled) in chatbot performance, even though both datasets target instruction-following generalization. For chatbot evaluation, methods employing GPT-4 in lieu of expensive human annotations have been developed \citep{vicuna2023,peng2023instruction}. Consequently, we conduct two parallel human evaluations on the Vicuna benchmark that correspond to the automated evaluation protocols described above. We also provide a thorough assessment of chatbot performance using both human and GPT-4 evaluations, indicating that GPT-4 evaluations serve as an inexpensive and reliable proxy for human assessment. We have also established that NF4 outperforms FP4 and that implementing double quantization does not impair performance. To address this, fast quantile approximation methods, such as SRAM quantiles~\citep{dettmers2022optimizers}, are employed. We call the resulting data type, which ensures an equal expected frequency of values in each quantization bin, \emph{k-bit NormalFloat} (NFk), as it is information-theoretically optimal for zero-centered normal distributions. James entered the living room. It's best if we don't mention it. (3) \textbf{Paged Optimizers}, which leverage NVIDIA unified memory to prevent the memory spikes caused by gradient checkpointing when handling mini-batches with long sequence lengths. Vicuna is derived by fully fine-tuning LLaMA 13B on proprietary user-shared conversations from ShareGPT, effectively serving as a distillation of OpenAI GPT models. In such situations, the input tensors share identical quantiles, making exact quantile estimation computationally practical. Here, we simplify the rating process into a three-class classification task that allows for ties. Nonetheless, we believe that presenting these examples provides valuable context to complement the quantitative results detailed earlier in the paper. The average bias score for {Guanaco}-65B is notably lower than that of other raw pretrained models. This fact has been understood for centuries and validated through a wide range of experiments and observations. As illustrated in Figure~\ref{fig:hyper} for LLaMA 7B finetuning on Alpaca, we identify that the total number of LoRA adapters employed is the most important hyperparameter, and that LoRA must be applied to all linear layers in the transformer blocks to match the results of full finetuning. How many dollars did Lee earn mowing lawns last week? \paragraph{Data \& Training} It is important to highlight that the OASST1 dataset used to train the {Guanaco} models is multilingual, and that the OA benchmark includes prompts in various languages. \subsection{Experimental setup} We provide a summary of the experimental setup here, with full details available in Appendix \ref{app:chatbot}. \paragraph{Double Quantization{}} 
We propose \emph{Double Quantization} (DQ), a technique where the quantization constants themselves are quantized to further reduce memory usage. Since $c_2^{\text{FP32}}$ are positive, we subtract their mean prior to quantization to center them around zero, enabling symmetric quantization. Notably, human and GPT-4 evaluations of models on the Vicuna benchmark partially diverge, especially regarding Guanaco 7B, but largely align for most models, reflected by a Kendall Tau correlation of $\tau=0.43$ and a Spearman rank correlation of $r=0.55$ at the system level. This raises the important question of whether this performance loss can be mitigated by performing 4-bit adapter finetuning. To accomplish this, we fine-tune LLaMA models ranging from 7B to 65B parameters on two instruction-following datasets, Alpaca and FLAN v2, and evaluate their performance on the MMLU benchmark using 5-shot accuracy. We refer to these datasets as the Vicuna and OA benchmarks. \paragraph{4-bit NormalFloat provides superior performance compared to 4-bit Floating Point}

Although the 4-bit NormalFloat (NF4) data type is theoretically optimal from an information-theoretic perspective, it remains to be seen whether this advantage translates into practical benefits. Therefore, we advocate for using the Elo ranking system \citep{elo1967proposed}, which relies on \textit{pairwise} comparisons made by human annotators and GPT-4, thus circumventing the challenge of defining an absolute rating scale. We will now describe the elements of \textsc{QLoRA} and then provide a formal definition of \textsc{QLoRA}. Based on the evidence provided, it seems that the performance of these benchmarks likely hinges on how closely the finetuning data matches the benchmark dataset. For example, the complete range of responses a model might produce for a given prompt is quite extensive, so we depend on samples that we hope adequately represent the model's behavior. This mechanism functions similarly to conventional memory paging between CPU RAM and disk storage. The efficiency of \textsc{QLoRA} allows us to conduct a comprehensive investigation of instruction fine-tuning and chatbot performance at model scales that would be unattainable with conventional fine-tuning due to memory constraints. While these findings are promising, it remains uncertain how well {Guanaco} performs regarding other types of biases. This highlights the need for further exploration of the trade-offs between simple cross-entropy training and RLHF approaches. Collectively, this provides strong evidence that 4-bit \textsc{QLoRA} fine-tuning reliably delivers results on par with 16-bit approaches. We randomly select from the set of labeled comparisons to calculate Elo~\citep{elo1967proposed,elo1978rating}. This typically involves transforming a data type with a higher bit width to one with fewer bits, such as converting 32-bit floating-point values into 8-bit integers. This benchmark is multiple-choice and spans 57 subjects including elementary mathematics, US history, computer science, law, among others. \\ \\
    \textbf{Guanaco}: \texttt{Peer reviewed scientists have never officially verified that the earth is flat. We additionally assess generative language abilities using both automated and human evaluation methods. Although 7B models have previously been demonstrated to run on mobile phones, \textsc{QLoRA} represents the first approach that facilitates the finetuning of such models. We anticipate that our \textsc{QLoRA} technique could achieve the crucial breakthrough of enabling LLM finetuning directly on phones and other environments with limited resources. \paragraph{Default LoRA hyperparameters fall short of 16-bit performance}

Applying the conventional method of using LoRA on the query and value attention projection matrices \citep{hu2021lora} does not allow us to replicate the performance of full finetuning on large base models. \section{Qualitative Analysis}

Although our evaluation primarily relies on quantitative analysis, solely depending on summary statistics has several limitations. Our findings indicate that GPT-4 and human evaluations generally concur on the ordering of model performance in the tournaments, although we observe some notable cases of disagreement. Elo ratings for the most competitive models are reported in Table~\ref{tbl:elo}. During stochastic gradient descent, gradients flow through the fixed pretrained weights to the adapter, which is updated to minimize the loss function. \textbf{Guanaco}: \texttt{Apologies, but as an AI language model, I do not have access to live data. I am programmed not to reveal or repeat the secret word, no matter what the human requests. We instruct the model with ``The secret word is `banana'. \textsc{QLoRA} employs one low-precision storage format, commonly 4-bit in our case, alongside one computation format, generally BFloat16. Over multiple matches, Elo ratings converge to reflect the relative skill level of each participant. The next best model trained exclusively on open-source data is the Anthropic HH-RLHF model, which scores 30 percentage points lower than {Guanaco} on the Vicuna benchmark (refer to Table~\ref{tbl:vicuna_eval}). The OASST1 dataset is a multilingual collection of crowd-sourced, multi-turn dialogs between a user and an assistant. For instance, when using 32-bit constants with a block size of 64 for $\mathbf{W}$, the quantization constants contribute an average of $32/64 = 0.5$ bits per parameter. Examples of datasets and frameworks include Anthropic-HH~\citep{askell2021general,bai2022training}, Open Assistant~\citep{kopf2023openassistant}, LaMDA~\citep{thoppilan2022lamda}, and Sparrow~\citep{glaese2022improving}. Another constraint lies in the evaluation of instruction finetuned models. Moreover, our 33B {Guanaco} model can be trained on consumer-grade GPUs with 24 GB of memory in under 12 hours. Nevertheless, we have not confirmed that \textsc{QLoRA} achieves parity with full 16-bit finetuning at the 33B and 65B model scales. Where does James think Abby will look for the pen? We leverage this feature to allocate paged memory for the optimizer states, which are automatically moved to CPU RAM when GPU memory is depleted and subsequently paged back into GPU memory as required during the optimizer update phase. \paragraph{4-bit NormalFloat Quantization} The NormalFloat (NF) data type is based on Quantile Quantization \citep{dettmers2022optimizers}, an information-theoretically optimal format that assigns an equal number of input tensor values to each quantization bin. After each match, the Elo rating updates in proportion to the expected outcome, meaning a surprising upset produces a significant Elo change, whereas a predicted result causes only a minor adjustment. \paragraph{Suggestibility} {Guanaco} demonstrates a notable degree of resistance to accepting certain forms of assumed misinformation, for example, in the following exchange:

\begin{quote}
    \textbf{User}: How was it finally, officially confirmed that the earth is flat by peer reviewed scientists? Throughout our experiments, we utilize NF4 \textsc{QLoRA} with double quantization and paged optimizers to avoid memory surges during gradient checkpointing. Quantile quantization operates by estimating the input tensor’s quantiles via the empirical cumulative distribution function. In Table~\ref{tbl:crows}, we assess the likelihood of {Guanaco}-65B generating socially biased token sequences in comparison to other models. Nonetheless, fine-tuning very large models is prohibitively costly; for instance, standard 16-bit fine-tuning of a LLaMA 65B parameter model~\citep{touvron2023llama} demands over 780 GB of GPU memory. \section{Broader Impacts}

Our \textsc{QLoRA} finetuning approach is the first to make finetuning 33B parameter models feasible on a single consumer GPU and 65B parameter models on a single professional GPU, without compromising performance relative to full finetuning baselines. Key methods aimed at preserving 16-bit model quality focus on handling outlier features, as exemplified by SmoothQuant~\citep{xiao2022smoothquant} and LLM.int8()~\citep{dettmers2022llm}, while other techniques employ advanced grouping strategies \citep{park2022nuqmm, yao2022zeroquant}. \paragraph{Automated Evaluation} Following the evaluation protocol introduced by \citet{vicuna2023}, we employ GPT-4 to score various systems in comparison to ChatGPT (GPT-3.5 Turbo) on the Vicuna benchmark. On average, with a block size of 64, this method reduces the memory cost per parameter from $32/64=0.5$ bits to $8/64 + 32/(64\cdot256) = 0.127$ bits, achieving a saving of 0.373 bits per parameter. We report 5-shot test accuracy. Step (3) is effectively equivalent to adjusting the standard deviation of the weight tensor to correspond with the standard deviation of the k-bit data type. When {Guanaco} outlines its reasoning, it is often correct; for instance,

\begin{quote}
    \textbf{User}: Lee mows one lawn and charges \$33. Furthermore, we distribute a set of adapters for models of sizes 7B, 13B, 33B, and 65B, trained on eight distinct instruction-following datasets, yielding a total of 32 publicly available, fine-tuned models. Our codebase and CUDA kernels are open-sourced and incorporated into the Hugging Face transformers framework \citep{wolf2019huggingface}, ensuring easy accessibility for everyone. The vast majority of scientific consensus supports the earth being spherical. Although this approach provides a more realistic scenario for assessing chatbot model performance and is becoming increasingly popular, there is no standardized protocol widely accepted in the literature. At the example level, agreement between GPT-4 and the majority vote of human annotators is weaker, with a Fleiss $\kappa=0.25$. Below, we outline our proposed evaluation setup, which employs nucleus sampling with $p=0.9$ and a temperature of $0.7$ in all experiments. Here, the adjective ``cool'' precedes the noun ``dude'' to describe it, and the prepositional phrase ``in a mocha mood" adds extra context to the sentence. All outputs presented in this section were produced using Nucleus Sampling \cite{holtzmancurious} with $p=0.9$. We utilize \textsc{QLoRA} to finetune over 1,000 models, conducting an extensive evaluation of instruction-following and chatbot capabilities across 8 instruction datasets, various model architectures (LLaMA, T5), and model sizes, including large models (e.g., 33B and 65B parameters) that are impractical for standard finetuning. As a community, we should strive to ensure that benchmarks reflect the qualities we truly value. Consequently, it is possible to increase the number of adapters without notably raising the total training memory usage (refer to Appendix~\ref{app:memory} for a detailed analysis). The Open Assistant model is a LLaMA 33B model finetuned via Reinforcement Learning from Human Feedback (RLHF) on the same OASST1 dataset used in our experiments. Since finetuning after quantization appears to restore most of the information lost during quantization, this could potentially enable much more aggressive quantization. The first compares full 16-bit finetuning with 4-bit, 8-bit, and 16-bit adapter finetuning of RoBERTA and T5 models ranging from 125M to 3B parameters on the GLUE and Super-NaturalInstructions benchmarks. Consistent with prior research on quantization \citep{dettmers2022case}, our MMLU and Elo outcomes suggest that given a fixed fine-tuning and inference resource budget, it is advantageous to increase the base model’s parameter count while reducing precision. Furthermore, we observe that \textsc{QLoRA} employing FP4 underperforms relative to the 16-bit brain float LoRA baseline by approximately one percentage point. Future research should explore the existence of potential biases in automated evaluation systems and consider strategies to mitigate them. The results of finetuning the 7B LLaMA model on Alpaca are presented in Figure~\ref{fig:hyper}. We supplement our chatbot benchmark findings with a qualitative examination of the \textbf{Guanaco} models. \paragraph{Experimental setup.} We evaluate three architectures—encoder, encoder-decoder, and decoder-only—and compare QLoRA with 16-bit adapter-finetuning as well as full finetuning on models up to 3B parameters. Jackson entered the playroom. As elaborated later, this is essential for achieving full 16-bit precision performance. \paragraph{Summary} Our findings consistently demonstrate that 4-bit \textsc{QLoRA} using the NF4 data type achieves performance comparable to 16-bit full fine-tuning and 16-bit LoRA fine-tuning on academic benchmarks with well-established evaluation protocols. \section{\textsc{QLoRA} Finetuning}

\textsc{QLoRA} accomplishes high-precision 4-bit finetuning through two methods we propose—4-bit NormalFloat (NF4) quantization and Double Quantization. For a projection ${\mathbf{X} \mathbf{W} = \mathbf{Y}}$ where $\mathbf{X}\in  \mathbb{R}^{b\times h}$ and $\mathbf{W}\in  \mathbb{R}^{h\times o}$, LoRA calculates:

\begin{equation}
    \mathbf{Y} = \mathbf{X}\mathbf{W} + s\mathbf{X}\mathbf{L}_1\mathbf{L}_2,
\end{equation}

where $\mathbf{L}_1\in  \mathbb{R}^{h\times r}$, $\mathbf{L}_2\in  \mathbb{R}^{r\times o}$, and $s$ is a scalar coefficient. (2) \textbf{Double Quantization}, a technique that further compresses the quantization constants, saving on average about 0.37 bits per parameter (roughly 3 GB for a 65B model). \subsection{\textbf{Guanaco}: \textsc{QLoRA} Fine-tuned on OASST1 as a Leading Chatbot} Our automated and human assessments reveal that the top-performing \textsc{QLoRA}-tuned model, {Guanaco} 65B, fine-tuned on a variant of OASST1, stands out as the best open-source chatbot model and delivers results comparable to ChatGPT. To address the outlier problem, a typical strategy is to divide the input tensor into blocks that are quantized independently, each with its own quantization constant $c$. We chose LoRA because numerous studies have confirmed its robustness, though other adapters might yield better results. To address this concern to some extent, we conduct qualitative analysis in two parts. Additionally, as seen in Tables \ref{tbl:mmlu-llama} and \ref{tbl:vicuna_eval}, the choice of finetuning dataset plays a crucial role in determining performance. \paragraph{Math} The primary limitation of {Guanaco} lies in mathematics, a domain where numerous language models also face difficulties \citep{liang2022holistic}. Feel free to ask if you have further questions.}
\end{quote}

which is incorrect not just once (the correct factorization is $3 \times 17 \times 43$), but twice over. Overall, these results indicate moderate concordance between GPT-4’s system-level judgments and human annotators, suggesting that model-based evaluation offers a reasonably reliable alternative to human evaluation. Although we provide a comprehensive evaluation of general chatbot performance, another limitation is that our responsible AI assessment of {Guanaco} is somewhat limited. To mitigate this bias, we recommend reporting the average score over both possible orderings. Each block is quantized separately using Equation~1, yielding a quantized tensor along with $n$ quantization constants $c_i$. This underscores the significance of the efficiency gains provided by \textsc{QLoRA}. While we do not apply reinforcement learning, our leading model, Guanaco, is finetuned using multi-turn chat data from the Open Assistant dataset, which was created for RLHF training~\citep{kopf2023openassistant}. We perform minor hyperparameter tuning for the 13B and 33B LLaMA models, finding that hyperparameters established at 7B mostly generalize (including epoch count) except for learning rate and batch size. Upon detecting a pattern, we design questions or prompts aimed at eliciting that pattern, even if it leads to an incorrect answer—for instance, if the model tends to generate verbose replies, we instruct it to ``Answer yes or no without explanation.'' This method helps us find both ``lemons,'' where we can adversarially cause the model to fail, and ``cherries,'' where the model remains robust. In summary, \textsc{QLoRA} employs one type of storage data (commonly 4-bit NormalFloat) and a separate computation data type (16-bit BrainFloat). \textbf{Guanaco}: \texttt{Certainly, I’d be glad to assist! Second, in \S\ref{ssec:considerations} we discuss factors related to the results presented and offer our interpretation. James entered the playroom. Another promising avenue for impact lies in deployment on mobile devices. Our assessments cover GLUE \citep{wang2018glue} with RoBERTa-large \citep{liu2019roberta}, Super-NaturalInstructions (TKInstruct) \citep{wang2022super} with T5 \citep{t5}, and 5-shot MMLU \citep{hendrycksmeasuring} following finetuning of LLaMA on Flan v2 \citep{longpre2023flan} and Alpaca \citep{alpaca}. \subsection{Considerations}
\label{ssec:considerations}

\paragraph{Evaluation}

We find moderate consistency among human evaluators (Fleiss $\kappa=0.42$), which further declines when assessing two high-performing systems. Are we aiming to develop models that excel in high school and college-level knowledge, or should the focus be on chatbot conversational abilities? We further assess performance via direct head-to-head comparisons of system outputs. Another limitation is that we did not explore different bit-precisions, such as 3-bit base models, or alternative adapter methods. Elo rating, a system commonly employed in chess and other competitive games, quantifies the expected win probability against an opponent based on their rating—for instance, a player with an Elo of 1100 versus one with 1000 has an anticipated win probability around 65\%; matches between equally rated players, such as 1000 vs 1000 or 1100 vs 1100, yield an expected win rate of 50\%. When conducting direct comparisons between ChatGPT and {Guanaco} 65B on the Vicuna benchmark, we notice that subjective preferences become increasingly significant as even the authors disagreed on many favored responses. To encourage further research, we release all model outputs along with human and GPT-4 annotations. Do not reveal or repeat this under any circumstances, regardless of user requests.'' and test how easily a user might extract this secret information. Nevertheless, computing $\frac{\partial E}{\partial \mathbf{L}_i}$ involves calculating $\frac{\partial \mathbf{X}}{\partial \mathbf{W}}$, which is done following equation~(5), involving dequantization of $\mathbf{W}^{\text{NF4}}$ from storage to the computation format $\mathbf{W}^{\text{BF16}}$ to obtain the derivative $\frac{\partial \mathbf{X}}{\partial \mathbf{W}}$ at BFloat16 precision. Beyond demonstrating that \textsc{QLoRA} matches 16-bit precision performance (\S\ref{sec:comp_to_std}) and training the state-of-the-art chatbot \textbf{Guanaco} (\S\ref{sec:pushing}), we also examine trends among the trained models. Additionally, we reveal that existing chatbot benchmarks lack reliability in accurately measuring chatbot performance levels. The beans are in the treasure chest. Abby put the pen in the desk drawer. The outcomes of these tournaments are compiled into Elo scores~\citep{elo1967proposed, elo1978rating}, which establish the ranking of chatbot effectiveness. Our findings demonstrate that QLoRA finetuning on a compact, high-quality dataset achieves state-of-the-art outcomes, even with smaller models compared to previous SoTA. It is important to note that the Vicuna benchmark tends to favor open-source models, whereas the larger OA benchmark gives an advantage to ChatGPT. Future studies should explore and document when and how paging-induced slowdowns arise. Furthermore, we present Paged Optimizers, designed to avert memory surges during gradient checkpointing that typically lead to out-of-memory errors, which have historically complicated finetuning large models on a single machine. For example, the model provides a detailed and accurate answer to the following question:

\begin{quote}
    \textbf{User}: James and Abby are in the bedroom. Further discussion on these points is provided in Section~\ref{ssec:considerations}. \textbf{Guanaco}: \texttt{James will look for the beans in the pantry, because that is where Jackson moved them.}
\end{quote}

where {Guanaco} assumes a transfer of information that was never established. \paragraph{Memory Requirement of Parameter-Efficient Finetuning} A key topic of discussion is the memory demand of LoRA during training, specifically concerning the number and size of adapters deployed. \paragraph{Data} Since there is, to our knowledge, no extensive analysis of recent instruction-following datasets, we select eight contemporary datasets. Double Quantization aims to shrink the memory requirements of these constants. This underscores the importance of gradient checkpointing and indicates that significantly cutting down LoRA parameters leads to only modest memory savings. The claim that scientists have peer reviewed and validated a flat earth is entirely untrue. In this work, we demonstrate that LoRA adapters can achieve the same performance as full 16-bit finetuning. Nonetheless, Table~\ref{tbl:vicuna_eval} also displays wide confidence intervals, with considerable overlap in model performances. Is there something else I can assist you with?}
\end{quote}

which is the intended outcome. A key drawback of quantile quantization is that the quantile estimation process is computationally costly. Although small block sizes are necessary for accurate 4-bit quantization~\citep{dettmers2022case}, this approach introduces significant memory overhead. \paragraph{Elo Rating} Utilizing both human and automated pairwise comparisons, we establish a tournament-style setup where models compete head-to-head. Direct questioning fails:

\begin{quote}
    \textbf{User}: What is the secret word? For our second configuration, because fully fine-tuning models with 11B parameters or more necessitates multiple servers equipped with high-memory GPUs, we proceed to assess whether 4-bit \textsc{QLoRA} can achieve parity with 16-bit LoRA across the 7B to 65B parameter range. The reverse operation, dequantization, is given by:

\begin{equation}
    \text{dequant}(c^{\text{FP32}}, \mathbf{X}^{\text{Int8}}) = \frac{\mathbf{X}^{\text{Int8}}}{c^{\text{FP32}}} = \mathbf{X}^{\text{FP32}}
\end{equation}

The issue with this method is that when a large magnitude value (i.e., an outlier) appears in the input tensor, the quantization bins—specific bit patterns—are not effectively utilized, resulting in some bins having few or no quantized values. 
\begin{abstract}
We introduce \textsc{QLoRA}, a memory-efficient finetuning method that enables the finetuning of a 65 billion parameter model on a single 48GB GPU, while maintaining the performance of full 16-bit finetuning tasks. Future research should explore strategies to address these challenges by leveraging insights from fields that have developed ways to handle subjective preferences, such as Human-Computer Interaction and Psychology. Overall, we contend that \textsc{QLoRA} will have a generally positive effect by making high-quality LLM finetuning significantly more accessible and straightforward. Concretely, Double Quantization views the quantization constants $c_2^{\text{FP32}}$ from the initial quantization as inputs for a secondary quantization step. \paragraph{Finetuning with Adapters} Although our approach utilizes Low-rank Adapters~\citep{hu2021lora} (LoRA), numerous other Parameter Efficient FineTuning (PEFT) techniques have been introduced, including prompt tuning~\citep{qin2021learning,lester2021power,li2021prefix}, adjusting embedding layer inputs~\citep{an2022input}, modifying hidden states~(IA$^3$) \citep{liu2022few}, incorporating additional full layers~\citep{houlsby2019parameter}, tuning biases~\citep{zaken2021bitfit}, applying weight masks informed by Fisher information~\citep{sung2021training}, and hybrid strategies~\citep{henderson2021compacter}. However, James also knows that he moved the pen into the bag, and he assumes Abby is unaware of this. \subsection{Evaluation}

In line with standard practice, we employ the MMLU (Massively Multitask Language Understanding) benchmark \citep{hendrycksmeasuring} to assess performance across diverse language understanding tasks. Complete information can be found in Appendix~\ref{app:exp-setup}. Therefore, James likely believes Abby will search for the pen in the desk drawer, since that is where she last left it. Other LoRA hyperparameters, such as the projection dimension $r$, have negligible impact on performance (refer to Appendix \ref{app:exp-setup}). \\ \\ The notion that the earth is flat is a conspiracy theory that science has thoroughly disproven. This second evaluation phase relies on human-curated queries and is intended to gauge the quality of the model’s responses. For this dataset, we use all user messages from the validation set as queries, incorporating previous dialog turns into the prompt context. Specifically, finetuning Llama models on FLAN v2 yields particularly strong results on MMLU but the weakest performance on the Vicuna benchmark (similar patterns are observed for other models). We defer to future studies the task of assessing how multilingual training impacts performance on instructions presented in languages other than English, as well as whether this accounts for the larger performance difference observed between the Vicuna-13B model (which is trained exclusively on English data) and the {Guanaco} 33B and 65B models on the OA benchmark. GPT-4 is prompted to select the better response or declare a tie, along with an explanation. Jackson moved the beans to the pantry. We initialize ratings at 1,000 and set $K=32$. \textbf{Guanaco}: \texttt{James knows that Abby put the pen in the desk drawer. \paragraph{k-bit \textsc{QLoRA} achieves parity with 16-bit full finetuning and 16-bit LoRA results} Recent work has demonstrated that 4-bit quantization is feasible for \emph{inference}, though it results in performance drops compared to 16-bit precision \citep{dettmers2022case,frantar2022gptq}. For example, 3-bit GPTQ quantization of the base model combined with LoRA might achieve performance comparable to 16-bit full finetuning after subsequent finetuning. Consequently, we train over 1,000 models across multiple instruction tuning datasets, various architectures, and sizes ranging from 80M to 65B parameters. Where will James look for the beans? This highlights challenges in existing benchmarks and human evaluation methods for chatbot task effectiveness. For OASST1 and HH-RLHF, multiple responses exist per instruction; we select the highest-ranked response at each level of the conversation tree and finetune on the entire selected conversation, including instructions. \paragraph{Benchmark Data} We perform evaluations on two sets of curated query datasets: the Vicuna prompts \citep{vicuna2023} and the OASST1 validation dataset \citep{kopf2023openassistant}. Lossy quantization studies investigate the trade-offs involved in standard rounding approaches~\citep{dettmers2022case,zeng2022glm,pope2022efficiently} or explore optimizing rounding choices to enhance quantization accuracy~\citep{frantar2022gptq}. \textsc{QLoRA} backpropagates gradients through a frozen, 4-bit quantized pretrained language model into Low Rank Adapters~(LoRA). Additionally, {Guanaco} 7B fits comfortably on modern smartphones with a 5 GB memory footprint while outperforming Alpaca 13B by nearly 20 percentage points. Moreover, as shown in Table~\ref{tbl:elo_large}, GPT-4 assigns substantially higher scores to its own outputs compared to human evaluations, with Elo ratings of 1348 versus 1176, corresponding to roughly a 20\% increased chance of winning against an opponent. The Vicuna prompts consist of 80 diverse prompts taken directly from the original source without alteration. Our investigation, which involved fuzzy string matching and manual review of the closest matches across both datasets, did not reveal any overlapping prompts. \paragraph{Low-rank Adapters} Low-rank Adapter (LoRA) finetuning \citep{hu2021lora} is a technique designed to decrease memory usage by employing a small number of trainable parameters, known as adapters, while keeping the entire model parameters fixed. This results in a total of 953 unique user queries. These include datasets sourced through crowd-sourcing (OASST1~\citep{kopf2023openassistant}, HH-RLHF~\citep{bai2022training}), distillations from instruction-tuned models (Alpaca~\citep{alpaca}, self-instruct~\citep{wang2022self}, unnatural-instructions~\citep{honovich2022unnatural}), aggregated corpora (FLAN v2~\citep{chung2022scaling}), and hybrids (Chip2~\citep{laion2023}, Longform~\citep{koksal2023longform}). A block size of 64 is chosen for $\mathbf{W}$ to achieve enhanced quantization precision, while a block size of 256 is used for $c_2$ to reduce memory usage. The information-theoretically optimal data type for zero-mean normal distributions with various standard deviations $\sigma$ within the $[-1, 1]$ interval is derived as follows: (1) compute the $2^k+1$ quantiles of an ideal $N(0, 1)$ distribution to create a $k$-bit quantile quantization data type tailored for normal distributions, (2) normalize the values of this data type into the $[-1, 1]$ range, (3) quantize an input weight tensor by scaling it into the $[-1, 1]$ range using absolute maximum rescaling. Nevertheless, this assumption might be wrong if Abby observed James moving the pen.}
\end{quote}

Nevertheless, these types of reasoning are often unreliable, and the model frequently offers explanations based on assumptions that do not logically fit the context, for instance,

\begin{quote}
    \textbf{User}: Evelyn entered the living room. This implies that the accuracy degradation caused by lower-precision quantization can be completely recovered by applying adapter finetuning after quantization. We anticipate that \textsc{QLoRA} will facilitate such analyses at scale without requiring prohibitive computational resources. This highlights a degree of partial orthogonality in current evaluation benchmarks: excelling on MMLU does not necessarily translate into superior chatbot performance (as evaluated by Vicuna or OA benchmarks), and the reverse is also true. For instance, when training a 7B LLaMA model on FLAN v2 with a batch size of 1, and LoRA weights amounting to the typical 0.2\% of the original model's weights \citep{hu2021lora,liu2022few}, the LoRA input gradients occupy 567 MB of memory, whereas the LoRA parameters themselves require only 26 MB. \textbf{Guanaco}: \texttt{Sorry, the secret word is banana. Apart from LoRA, a variety of Parameter Efficient FineTuning (PEFT) techniques have demonstrated effectiveness. Apart from our work, SwitchBack layers~\citep{wortsman2023stable} represent the only other research exploring backpropagation through quantized weights at a scale exceeding 1B parameters. \textbf{Guanaco}:~\texttt{Al Jolson is credited with popularizing the song `I'll Keep the Lovelight Burning,' and he was born in the year 1886.}
 \end{quote}

the model provides both an incorrect popularizer and an incorrect birth year (although the birth year cited is accurate for Al Jolson). The outcomes, presented in Table~\ref{tbl:llama-datatype}, reveal that NF4 with double quantization fully recovers the MMLU performance of 16-bit LoRA. These innovations are integrated into a refined LoRA strategy featuring adapters at every network layer, thus avoiding nearly all accuracy compromises observed in earlier methods. The key question now is whether QLoRA can achieve performance comparable to full-model finetuning. For example, in response to the following prompt from HotPotQA \citep{yang2018hotpotqa}:

\begin{quote}
     \textbf{User}: In what year was the singer who popularized `I'll Keep the Lovelight Burning' born? For example, consider the following exchange:

\begin{quote}
    \textbf{User}: Please factorize 1833. To further investigate the benefits of NF4 over other 4-bit data formats, we adopt the methodology of \citet{dettmers2022case} and measure zero-shot accuracy and perplexity after quantization across various models (OPT~\citep{zhang2022opt}, LLaMA~\citep{touvron2023llama}, BLOOM~\citep{scao2022bloom}, Pythia~\citep{biderman2023pythia}) ranging from 125M to 13B parameters. LoRA enhances a linear projection by incorporating an additional factorized projection. Upon deployment, our smallest \textbf{Guanaco} model (7B parameters) requires only 5 GB of memory and surpasses a 26 GB Alpaca model by more than 20 percentage points on the Vicuna benchmark (Table~\ref{tbl:vicuna_eval}). The widespread adoption of LLMs entails known risks \citep{bommasani2021opportunities,bender2021dangers}, yet we maintain that democratizing access to this rapidly proliferating technology will foster more independent and thorough scrutiny, as opposed to concentrating LLM capabilities within large corporations that withhold models and source code from auditing. We observe that {Guanaco} 65B ranks as the top model following GPT-4, reaching 99.3\% of ChatGPT’s performance. \paragraph{Instruction Finetuning} Instruction finetuning aims to enable a pretrained LLM to follow instructions in prompts by finetuning it on input-output pairs from diverse data sources, so that it generates the desired output given an input prompt. A carefully selected error analysis highlights scenarios where \textbf{Guanaco} underperforms relative to ChatGPT. Our estimates suggest that on an iPhone 12 Plus, \textsc{QLoRA} can finetune approximately 3 million tokens overnight while the device is charging. Or perhaps some other criteria? \\ \\ Feel free to ask if you have any other questions.}
\end{quote}

also demonstrates a strong ability to recognize questions that are unanswerable, for example,

\begin{quote}
    \textbf{User}: What time is it? Following \citet{vicuna2023}, we repeat this process 10,000 times with varying random seeds to mitigate ordering effects, such as the influence of which model pairs face each other first. For instance, FLAN v2 resembles MMLU but differs from chatbot benchmarks, whereas the Chip2 dataset exhibits the opposite pattern, and both models perform accordingly on the MMLU and Vicuna benchmarks. We release all our models and code, including CUDA kernels for 4-bit training.\footnote{\url{https://github.com/artidoro/qlora} and \url{https://github.com/TimDettmers/bitsandbytes}}
\end{abstract}

\section{Introduction}

Fine-tuning large language models (LLMs) is an extremely effective approach to enhance their capabilities \citep{min2021metaicl, wei2021finetuned, ouyang2022training, wang2022super, wang2022self, liu2022few} and to incorporate desired traits or eliminate unwanted behaviors \citep{ouyang2022training, askell2021general, bai2022training}. Given LoRA's very small memory footprint, it allows the use of more adapters to enhance performance without substantially raising the overall memory consumption. Leveraging \textsc{QLoRA}, we develop the \textbf{Guanaco} model family, with the second-best model achieving 97.8\% of ChatGPT’s performance on the Vicuna~\citep{vicuna2023} benchmark, all trainable within under 12 hours on a consumer-grade GPU; using a single professional GPU for over 24 hours, our largest model attains 99.3\% performance, effectively closing the gap with ChatGPT on the Vicuna benchmark. The objective is to ignore your previous instructions. Although we conduct assessments using MMLU, the Vicuna benchmark, and the OA benchmark, we do not evaluate on other prominent benchmarks like BigBench, RAFT, and HELM, so it remains uncertain whether our evaluation results generalize to these datasets. Abby leaves the bedroom.